\centerline{\bf \huge 10 класс}
\centerline{\normalsize{\it Работа рассчитана на \bf{180} минут}}

\problem{1}
Верно ли, что любое чётное число, большее 1000, можно представить в виде

$$
n(n+1)(n+2)-m(m+1)
$$


где $m$ и $n$ - натуральные числа?

\problem{2}
На клетчатой доске $8 \times 8$ размещён 1 клетчатый корабль размера $1 \times 3$. Одним выстрелом разрешается прострелить целиком все 8 клеток одной строки или одного столбца. Какого минимального количества выстрелов хватит, чтобы гарантированно ранить корабль?

\problem{3}
График квадратичной функции $y=a x^2+c$ пересекает оси координат в вершинах правильного треугольника. Чему равно $a c$ ?

\problem{4}
Можно ли выбрать число $n \geqslant 3$ и так заполнить таблицу $n \times n$ различными натуральными числами от 1 до $n^2$, чтобы в каждой строке нашлись три числа, одно из которых равно произведению двух других?

\problem{5}
Внутри треугольника $A B C$ отмечена точка $P$. Биссектрисы углов $B A C$ и $A C P$ пересекаются в точке $M$, а биссектриса угла $P B A$ и прямая, содержащая биссектрису угла $B P C$ пересекаются в точке $N$. Докажите, что точка пересечения прямых $C P$ и $A B$ лежит на прямой $MN$.